\begin{longtable}{|p{3.3cm}|p{3.3cm}|p{3.3cm}|p{3.3cm}|}
\caption{So sánh ba hướng nền tảng chuỗi cho hệ thống văn bằng số}
\label{tab:blockchain-comparison} \\
\hline
\textbf{Tiêu chí} & \textbf{Chuỗi công khai hợp đồng chung} & \textbf{Chuỗi được cấp phép} & \textbf{Chuỗi công khai máy ảo Move} \\
\hline
\endfirsthead

\hline
\textbf{Tiêu chí} & \textbf{Chuỗi công khai hợp đồng chung} & \textbf{Chuỗi được cấp phép} & \textbf{Chuỗi công khai máy ảo Move} \\
\hline
\endhead

\hline
\endfoot

\hline
\endlastfoot

Ngôn ngữ hợp đồng
& Solidity
& Go, Java, JavaScript
& Move
\\ \hline

Mô hình lập trình
& Hướng tài khoản
& Hướng tài khoản
& Hướng tài nguyên, hệ thống kiểu ngăn nhân bản và hủy ngoài ý muốn
\\ \hline

Kiểm soát truy cập giữa các mô-đun
& Theo quy ước và bộ sửa đổi do lập trình viên viết
& Theo chính sách của mạng
& Cơ chế \texttt{friend} ở tầng ngôn ngữ
\\ \hline

Thời gian xác nhận
& Khoảng 12 giây
& Dưới 1 giây
& Dưới 1 giây, nhờ thực thi song song
\\ \hline

Đọc dữ liệu không cần cấp phép
& Có
& \textbf{Không}: bên ngoài tập thành viên phải đi qua một cổng dịch vụ
& Có
\\ \hline

Chi phí ở giai đoạn thử nghiệm
& Cao và biến động theo thị trường
& Cần vận hành nút xác thực riêng
& Thấp trên mạng thử nghiệm
\\ \hline

Phù hợp cho kiểm chứng xác minh độc lập
& Có
& Không: không kiểm chứng được tính chất {``}không cần thỏa thuận trước{''}
& Có
\\ \hline

Phù hợp cho triển khai thật nhiều cơ sở
& Chi phí là rào cản
& Phù hợp nếu theo mô hình liên hợp
& Cần đánh giá lại chi phí và quản trị
\\ \hline

\end{longtable}
